\documentclass[11pt]{article}
\usepackage{amsmath,amssymb,amsfonts}
\usepackage[margin=1in]{geometry}
\usepackage{enumitem}
\usepackage{booktabs}
\usepackage{hyperref}
\title{Harnessing Rydberg Atomic Receivers (v5, accepted) --- Complete Equation-by-Equation Reference}
\author{Compiled from the accepted v5 manuscript (arXiv:2501.11842)}
\date{}
\begin{document}
\maketitle
\tableofcontents

\section{The three receiver types --- overview}

The paper compares \textbf{three} receiver architectures throughout.
\begin{enumerate}
  \item \textbf{Conventional RF receiver}: classical dipole antenna + LNA
  superheterodyne chain. Pure baseline, no Rydberg physics.
  \item \textbf{LO-free Rydberg atomic receiver}: reads the RF field from
  the Autler--Townes (AT) splitting of the EIT peak directly. Measures
  amplitude only.
  \item \textbf{LO-dressed Rydberg atomic receiver}: a strong extra LO RF
  tone is injected on-resonance with the same Rydberg transition, turning
  the atoms into a coherent mixer. Measures amplitude \emph{and} phase,
  and is the most sensitive of the three.
\end{enumerate}

\begin{center}
\begin{tabular}{@{}ll@{}}
\toprule
Figure & What it shows / which receiver(s) \\
\midrule
Fig.~1 & Schematic only (energy levels, LO-free vs LO-dressed setups) \\
Fig.~2 & Noise-chain schematic, all three receivers (block diagram only) \\
Fig.~3 & $P_\text{out}$ vs $\Delta_c,\Omega_\text{RF}$ --- LO-free, no link budget yet \\
Fig.~4 & $P_\text{out}$ heatmap vs $\Delta_c,\mathcal{R}$ --- LO-free resolvability \\
Fig.~5 & $P_\text{out}$ vs $\Delta_c, d_\text{Tx-Rx}$ + SNR --- LO-free, full link budget \\
Fig.~6 & Linear dynamic range (LDR) of the LO-dressed receiver \\
Fig.~7 & LDR translated into $P_\text{Rx}$ bounds vs $\Omega_\text{LO}$ --- LO-dressed \\
Fig.~8 & SNR vs distance --- all three receivers compared \\
Fig.~9 & Mutual information vs $P_\text{Rx}$ --- all three receivers \\
Fig.~10 & SER vs $P_\text{Rx}$ --- all three receivers \\
\bottomrule
\end{tabular}
\end{center}

\section{Section II: Preliminaries}

\subsection*{Eq.~(1) --- the AT-splitting readout relation}
\begin{equation}
\Delta\nu =
\begin{cases}
\dfrac{\Omega_\text{RF}}{2\pi}, & \text{scanning the coupling laser},\\[2mm]
\dfrac{\Omega_\text{RF}}{2\pi}\dfrac{f_p}{f_c}, & \text{scanning the probe laser}.
\end{cases}
\tag{1}
\end{equation}
\textbf{Terms:} $\Delta\nu$ is the frequency separation between the two
split EIT peaks that you actually observe on the photodetector.
$\Omega_\text{RF}$ is the RF Rabi frequency (what you are trying to
measure). $f_p,f_c$ are the probe and coupling laser optical
frequencies.
\textbf{Significance:} this is the foundational readout equation of the
whole LO-free concept: it says the splitting you literally see encodes
$\Omega_\text{RF}$ directly. Scanning the coupling laser gives the
clean, calibration-free form $\Delta\nu=\Omega_\text{RF}/2\pi$ (no probe
absorption background), which is why it's preferred experimentally.
\textbf{Not used for Fig.~3 itself} (Fig.~3 plots $P_\text{out}$ directly
against $\Delta_c$, not $\Delta\nu$) --- but it is the conceptual
justification for why Fig.~3's $x$-axis is $\Delta_c$ at all.

\subsection*{Eq.~(2) --- Beer--Lambert probe transmission}
\begin{equation}
P_\text{out}(t) = P_\text{in}\exp\{-k_pL\,\Im(\chi(t))\}. \tag{2}
\end{equation}
\textbf{Terms:} $P_\text{out}(t)$ = transmitted probe laser intensity
measured on the photodetector --- \emph{this is the quantity plotted on
the $z$-axis of Fig.~3, and of every other $P_\text{out}$ figure in the
paper}. $P_\text{in}$ = incident probe laser intensity (before the vapor
cell). $k_p=2\pi/\lambda_p$ = probe wavenumber. $L$ = vapor-cell length.
$\chi(t)$ = complex electric susceptibility of the atomic medium at time
$t$; only its imaginary part matters here because $\Im(\chi)$ is what
controls absorption (the real part would control dispersion/phase
shift, not amplitude loss).
\textbf{Significance:} this is a literal, standard Beer--Lambert
absorption law, with the medium's absorption coefficient written as
$k_p\Im(\chi)$. \textbf{This is the single formula that turns the
quantum-mechanical coherence into the number plotted in Fig.~3} ---
every value in our 3D surface is computed by plugging a value of
$\Im(\chi)$ into this exact exponential. If Fig.~3 doesn't match the
paper, the error must live either (a) upstream, in what $\chi(t)$ or
$\Im(\chi)$ actually evaluates to (wrong Hamiltonian, wrong steady
state, wrong parameters), or (b) in the constants $k_p$, $L$, $P_\text{in}$
fed into this exact exponential. There is no other place for a bug to
hide in this equation itself --- it's a one-line formula.

\subsection*{Eq.~(3) --- $P_\text{in}$ in terms of the probe Rabi frequency}
\begin{equation}
P_\text{in} = \frac{\pi}{2Z_0}\left(\frac{d\,\Omega_p\hbar}{2\sqrt{\wp_{12}}}\right)^2. \tag{3}
\end{equation}
\textbf{Terms:} $d$ = probe beam $1/e^2$ diameter. $\Omega_p$ = probe
Rabi frequency. $\hbar$ = reduced Planck's constant. $\wp_{12}$ = dipole
moment of the $|1\rangle$--$|2\rangle$ transition (\textbf{note}: as
written, $\wp_{12}$ is used \emph{under a square root} here --- confirms
$\wp_{12}$ is defined as an already-squared quantity, i.e.
$\wp_{12}=(2.5ea_0)^2$ as stated later in the paper, not
$\wp_{12}=2.5ea_0$). $Z_0$ = free-space impedance.
\textbf{Significance:} this equation is the physical link between the
probe beam's intensity/power and its Rabi frequency (via the beam
diameter $d$, which sets the field's spatial concentration). \textbf{This
is the exact equation at the center of the known Pin/diameter
inconsistency}: plugging in the paper's own stated $d=0.76$~mm and
$\Omega_p/2\pi=8$~MHz gives $P_\text{in}=39.08\,\mu$W, not the
$20.7\,\mu$W the paper also states directly in the same parameter list.
\textbf{This equation is not directly used to compute Fig.~3's
$P_\text{out}$ surface} (we use the stated $P_\text{in}=20.7\,\mu$W
directly in Eq.~2, and the stated $\Omega_p/2\pi=8$~MHz directly in the
Hamiltonian, Eq.~6) --- but it is the reason those two paper-stated
numbers can never both be exactly self-consistent, which is one
confirmed, real source of an overall multiplicative scale mismatch
(though NOT of any shape/sharpness mismatch, since Eq.~3 only sets an
overall constant, not how $P_\text{out}$ depends on $\Delta_c$ or
$\Omega_\text{RF}$).

\subsection*{Eq.~(4) --- susceptibility in terms of the coherence}
\begin{equation}
\chi(t) = C_0\,\rho_{21}(t), \qquad C_0 = \frac{-2N_0\wp_{12}}{\epsilon_0\hbar\Omega_p}. \tag{4}
\end{equation}
\textbf{Terms:} $\rho_{21}(t)$ = the $|1\rangle$--$|2\rangle$ coherence,
an off-diagonal element of the atomic density matrix --- \emph{the
single quantity that the entire quantum-mechanical part of the model
(Hamiltonian + Lindblad + steady state) exists to compute}. $C_0$ = a
real, negative constant bundling $N_0$ (atomic density), $\wp_{12}$
(dipole moment), $\epsilon_0$ (vacuum permittivity), $\hbar$, and
$\Omega_p$. $N_0$ = total atomic number density in the vapor cell.
\textbf{Significance:} this equation is the direct bridge between
``quantum physics'' ($\rho_{21}$, computed via QuTiP) and ``classical
optics'' ($\chi$, fed into Beer--Lambert, Eq.~2). \textbf{$C_0$ is
exactly where $N_0=4.89\times10^{16}\,\text{m}^{-3}$ enters the whole
calculation} --- since $C_0$ is a simple linear prefactor, doubling
$N_0$ exactly doubles $\Im(\chi)$ for the same $\rho_{21}$, which (through
Eq.~2's exponential) can massively change $P_\text{out}$'s absolute
scale and apparent sharpness. This is the equation responsible for the
already-confirmed ``high-$N_0$ compression'' effect: because $N_0$ is
large, small physically-correct changes in $\rho_{21}$ (from changing
$\Delta_c$ or $\Omega_\text{RF}$ by a little) get amplified into large
changes in $\Im(\chi)$, and then further amplified exponentially by
Eq.~2 --- producing a sharper, narrower-looking $P_\text{out}$ curve than
a naive linear intuition would suggest.

\subsection*{Eq.~(5) --- the master equation}
\begin{equation}
\dot{\boldsymbol\rho} = \frac{\partial\boldsymbol\rho}{\partial t} = -\frac{j}{\hbar}[\mathbf H,\boldsymbol\rho] + \mathcal L. \tag{5}
\end{equation}
\textbf{Terms:} $\boldsymbol\rho$ = the full $4\times4$ density matrix of
the 4-level atomic system. $\mathbf H$ = the system Hamiltonian
(Eq.~6). $[\mathbf H,\boldsymbol\rho]=\mathbf H\boldsymbol\rho-\boldsymbol\rho\mathbf H$
= the commutator, generating coherent (unitary) evolution. $\mathcal L$
= the Lindblad dissipator (Eq.~7), generating incoherent
(spontaneous-emission) evolution. $j=\sqrt{-1}$ (engineering notation
for the imaginary unit, used throughout the paper instead of $i$).
\textbf{Significance:} this is the standard Lindblad master equation for
an open quantum system --- the general-purpose equation that
\texttt{qutip.steadystate()} solves (by construction) when you pass it
$\mathbf H$ and the collapse operators. \textbf{There is no ambiguity or
possible bug in this equation itself} --- QuTiP's steady-state solver
implements exactly this, given a correct $\mathbf H$ and a correct set
of collapse operators. Any discrepancy must trace back to $\mathbf H$
(Eq.~6) or the collapse operators (Eq.~7), not to this equation.

\subsection*{Eq.~(6) --- the 4-level Hamiltonian}
\begin{equation}
\mathbf H = \frac{\hbar}{2}
\begin{bmatrix}
0 & \Omega_p & 0 & 0\\
\Omega_p & -2\Delta_p & \Omega_c & 0\\
0 & \Omega_c & -2(\Delta_p+\Delta_c) & \Omega\\
0 & 0 & \Omega & -2(\Delta_p+\Delta_c+\Delta_\text{RF})
\end{bmatrix}. \tag{6}
\end{equation}
\textbf{Terms:} basis states $|1\rangle,|2\rangle,|3\rangle,|4\rangle$
correspond to Cs levels $6S_{1/2}, 6P_{3/2}, 47D_{5/2}, 48P_{3/2}$.
$\Omega_p$ = probe Rabi frequency (couples $|1\rangle\leftrightarrow|2\rangle$).
$\Omega_c$ = coupling Rabi frequency (couples
$|2\rangle\leftrightarrow|3\rangle$). $\Omega$ (bare, unsubscripted) =
the $|3\rangle\leftrightarrow|4\rangle$ coupling --- \textbf{the paper
explicitly states in the text immediately after this equation that
``the Rabi frequency $\Omega$ in (6) differs for LO-free and LO-dressed
Rydberg atomic receivers.''} For the LO-free receiver (Fig.~3, 4, 5),
$\Omega=\Omega_\text{RF}$ directly (confirmed by Eq.~8's own use of
$\Omega_\text{RF}$ in exactly this matrix slot). For the LO-dressed
receiver (Fig.~6, 7), $\Omega=\Omega_\text{total}=\Omega_\text{LO}+e^{j(2\pi\Delta f\,t+\Delta\phi)}\Omega_\text{RF}$
(a complex phasor, defined later at Eq.~26/28 context). $\Delta_p,\Delta_c,\Delta_\text{RF}$
= detunings of the probe, coupling, and RF fields from their respective
transitions.
\textbf{Significance:} this is \emph{the} Hamiltonian used for every
single QuTiP steady-state solve in this whole reproduction (Fig.~3
through Fig.~7 all reuse this exact 4-level structure, changing only
which detunings are swept and what value goes in the $\Omega$ slot).
\textbf{For debugging Fig.~3 specifically}: confirm your code's matrix
uses $-\Delta_p$ (not $-2\Delta_p$) on the diagonal if your Hamiltonian
is defined directly in angular-frequency units without an explicit
$\hbar/2$ prefactor (i.e., if you write \texttt{H = ...} without
dividing by $\hbar$ afterward) --- dividing the paper's $\frac{\hbar}{2}\times(-2\Delta_p)$
by $\hbar$ gives exactly $-\Delta_p$, which is what a
Hamiltonian-in-frequency-units convention (the usual QuTiP convention
for a Hermitian generator with $\hbar=1$) should use. This has already
been verified correct in this project's code.

\subsection*{Eq.~(7) --- the Lindblad dissipator}
\begin{equation}
\mathcal L =
\begin{bmatrix}
\gamma_2\rho_{22} & -\gamma_{12}\rho_{12} & -\gamma_{13}\rho_{13} & -\gamma_{14}\rho_{14}\\
-\gamma_{21}\rho_{21} & \gamma_3\rho_{33}-\gamma_2\rho_{22} & -\gamma_{23}\rho_{23} & -\gamma_{24}\rho_{24}\\
-\gamma_{31}\rho_{31} & -\gamma_{32}\rho_{32} & \gamma_4\rho_{44}-\gamma_3\rho_{33} & -\gamma_{34}\rho_{34}\\
-\gamma_{41}\rho_{41} & -\gamma_{42}\rho_{42} & -\gamma_{43}\rho_{43} & -\gamma_4\rho_{44}
\end{bmatrix},\qquad \gamma_{ij}=\frac{\gamma_i+\gamma_j}{2}. \tag{7}
\end{equation}
\textbf{Terms:} $\gamma_2,\gamma_3,\gamma_4$ = spontaneous-emission
(population decay) rates out of levels 2, 3, 4 respectively ($\gamma_1=0$
since $|1\rangle$ is the ground state and cannot decay further).
$\gamma_{ij}=(\gamma_i+\gamma_j)/2$ = the dephasing rate of coherence
$\rho_{ij}$, a standard relation for decay-induced dephasing.
\textbf{Significance:} this matrix describes a simple \emph{cascade
decay} chain: $|2\rangle\to|1\rangle$ at rate $\gamma_2$,
$|3\rangle\to|2\rangle$ at rate $\gamma_3$, $|4\rangle\to|3\rangle$ at
rate $\gamma_4$ --- population only ever flows \emph{downward}. This
exact structure is generated automatically by three QuTiP collapse
operators $\sqrt{\gamma_2}|1\rangle\langle2|$,
$\sqrt{\gamma_3}|2\rangle\langle3|$, $\sqrt{\gamma_4}|3\rangle\langle4|$
--- if your code's collapse operators are anything other than exactly
these three (e.g., missing one, wrong ket/bra order, or an extra
spurious dephasing-only operator), the resulting steady state will
differ from the paper's model, silently. \textbf{This has already been
verified correct} in this project's Fig.~3 code by direct comparison of
the collapse-operator list against this matrix structure.

\subsection*{Eq.~(8) --- closed-form steady-state coherence (weak-probe limit)}
\begin{equation}
\rho_{21} = \cfrac{-j(\Omega_p/2)}{\gamma_{21}-j\Delta_p-\cfrac{(\Omega_c/2)^2}{\gamma_{31}-j(\Delta_p+\Delta_c)-\cfrac{(\Omega_\text{RF}/2)^2}{\gamma_{41}-j(\Delta_p+\Delta_c+\Delta_\text{RF})}}}. \tag{8}
\end{equation}
\textbf{Terms:} a nested continued fraction in $\Omega_p,\Omega_c,\Omega_\text{RF}$
and the three detunings, using the same $\gamma_{ij}$ dephasing rates as
Eq.~7.
\textbf{Significance:} this is the paper's own \emph{analytic}
alternative to a full QuTiP steady-state solve, derived (Appendix A)
under the weak-probe approximation $\rho_{11}\approx1$,
$\rho_{22},\rho_{33},\rho_{44}\approx0$. \textbf{Verified this session
to be numerically pathological in Fig.~3's actual operating regime}: at
$\Delta_c=0$ and $\Omega_\text{RF}=0$ it gives unphysical gain
($P_\text{out}>P_\text{in}$), and for any $\Omega_\text{RF}\gtrsim0.5$~MHz
it collapses to numerically zero and stays flat --- because
$\gamma_4/2\pi=1.7$~kHz is so tiny that the innermost term
$(\Omega_\text{RF}/2)^2/\gamma_{41}$ explodes almost immediately. Tried
three different sign conventions (flip the outer $j$, flip all $j$'s,
flip only the detuning signs) --- none produce a sane curve. \textbf{Do
not use Eq.~8 to try to reproduce or debug Fig.~3} --- the paper's own
Fig.~3 was generated with full QuTiP (stated directly in Sec.~V-A: ``We
employ the QuTiP toolkit for simulating the quantum system''), not this
closed form. This equation is only useful for physical intuition and
for deriving other analytic results elsewhere in the paper (e.g., the
LO-dressed Appendix~C derivations).

\subsection*{Eq.~(9) --- Maxwell--Boltzmann velocity distribution}
\begin{equation}
g(v_z) = \frac{1}{\sqrt{2\pi}\sigma_v}e^{-v_z^2/(2\sigma_v^2)}. \tag{9}
\end{equation}
\textbf{Terms:} $v_z$ = atomic velocity component along the laser
propagation axis. $\sigma_v=\sqrt{k_BT/m_\text{Rb}}$ = the thermal
velocity spread (standard deviation), with $k_B$ = Boltzmann constant,
$T$ = vapor temperature, $m_\text{Rb}$ = atomic mass.
\textbf{Significance:} standard 1-D thermal velocity distribution for a
room-temperature atomic vapor --- the starting point for modeling
Doppler broadening.

\subsection*{Eq.~(10) --- Doppler-shifted detunings}
\begin{equation}
\Delta_p(v_z)=\Delta_p-\frac{2\pi}{\lambda_p}v_z,\qquad
\Delta_c(v_z)=\Delta_c+\frac{2\pi}{\lambda_c}v_z. \tag{10}
\end{equation}
\textbf{Terms:} $\lambda_p,\lambda_c$ = probe/coupling laser
wavelengths. The opposite signs on the two Doppler-shift terms come
directly from the probe and coupling lasers being
\textbf{counter-propagating} (stated explicitly in the text) --- an
atom moving toward one beam is moving away from the other.
\textbf{Significance:} this is how a single atomic velocity $v_z$ maps
onto an \emph{effective}, atom-specific pair of detunings that then get
substituted into Eq.~8 (or the full Hamiltonian) before averaging.
\textbf{Directly relevant to the Fig.~3 debugging question}: this
project deliberately does \emph{not} apply Eq.~10/11 to Fig.~3, because
a real convergence test showed that properly-converged Doppler
averaging completely erases the AT splitting (a $\sim$0\% dip at
$\Delta_c=0$), which contradicts the real published Fig.~3 (which
clearly shows resolved splitting). The paper's own cited source for
this exact kind of EIT-AT spectrum plot explicitly states it uses a
\textbf{cold-atom model} (no Doppler averaging) for its own analogous
figure. If you want to re-litigate this decision as part of your deep
dive, this is the equation pair to revisit --- but the convergence-test
evidence against including it is strong.

\subsection*{Eq.~(11) --- Doppler-averaged coherence}
\begin{equation}
\bar\rho_{21}(\Delta_p)=\int_{-\infty}^{+\infty}\rho_{21}(v_z)\,g(v_z)\,dv_z. \tag{11}
\end{equation}
\textbf{Terms:} $\rho_{21}(v_z)$ = the coherence computed (via Eq.~8, or
a full solve) using the velocity-shifted detunings from Eq.~10, for one
specific atomic velocity class $v_z$. The integral averages over the
full thermal velocity distribution $g(v_z)$ from Eq.~9.
\textbf{Significance:} this is the final Doppler-averaging step ---
in principle, the ``correct'' $\rho_{21}$ to feed into Eq.~4/2 for a
warm vapor. As noted under Eq.~10, this project's convergence testing
showed this integral needs $\gtrsim400$ velocity-class samples to
converge (given the single-photon Doppler width is $\sim$158~MHz versus
the $\sim$MHz-scale EIT features), and a properly-converged version
erases the AT splitting entirely --- which is why Fig.~3 in this
project uses the \emph{unaveraged}, single-velocity-class ($v_z=0$,
i.e. cold-atom) result instead.

\section{Section III-A: LO-free Rydberg atomic receiver model}

\subsection*{Eq.~(12) --- RF Rabi frequency from the field}
\begin{equation}
\Omega_\text{RF} = \frac{|E_\text{RF}|\wp_\text{RF}}{\hbar}. \tag{12}
\end{equation}
\textbf{Terms:} $E_\text{RF}$ = incident RF electric field amplitude.
$\wp_\text{RF}$ = dipole moment of the $|3\rangle$--$|4\rangle$
(Rydberg--Rydberg) transition, stated as $\wp_\text{RF}=-1443.459\,ea_0$.
\textbf{Significance:} the fundamental field-to-Rabi-frequency
conversion, used any time a distance or transmit power needs to become
an $\Omega_\text{RF}$ value to feed into the quantum model (Fig.~5,
Fig.~7). Not used by Fig.~3 itself (Fig.~3 sweeps $\Omega_\text{RF}$
directly as an independent axis, with no link budget involved).

\subsection*{Eq.~(13) --- LO-free signal model and received power definition}
\begin{equation}
y = \sqrt{P_\text{Rx}}\,hx+n_\text{ex}, \qquad P_\text{Rx}=\frac{|E_\text{RF}|^2}{2Z_0}. \tag{13}
\end{equation}
\textbf{Terms:} $x\sim\mathcal{CN}(0,1)$ = unit-power complex Gaussian
transmitted symbol. $h\sim\mathcal{CN}(0,1)$ = complex Gaussian channel
gain. $n_\text{ex}$ = extrinsic (BBR) noise. $P_\text{Rx}$ = received RF
power, defined here directly from the field amplitude.
\textbf{Significance:} sets up the standard wireless-communications
signal model wrapping around the Rydberg physics --- not used for
Fig.~3 (a pure physics figure with no wireless-link elements).

\subsection*{Eq.~(14) --- Friis transmission equation}
\begin{equation}
P_\text{Rx} = P_\text{Tx}G_\text{Tx}G_\text{Rx}\left(\frac{\lambda}{4\pi d_\text{Tx-Rx}}\right)^2. \tag{14}
\end{equation}
\textbf{Terms:} $P_\text{Tx},G_\text{Tx}$ = transmit power and Tx
antenna gain. $G_\text{Rx}$ = classical receive antenna gain.
$d_\text{Tx-Rx}$ = link distance. $\lambda$ = RF wavelength.
\textbf{Significance:} the standard classical link-budget formula,
quoted here explicitly as \emph{not directly applicable} to the Rydberg
receiver, since ``there is no metallic antenna and hence no meaningful
classical gain $G_\text{Rx}$'' --- this sets up why Eq.~15--17 are needed
instead.

\subsection*{Eq.~(15) --- received power flux density}
\begin{equation}
S_\text{Rx} = \frac{P_\text{Tx}G_\text{Tx}}{4\pi d_\text{Tx-Rx}^2}. \tag{15}
\end{equation}
\textbf{Terms:} $S_\text{Rx}$ = power flux density (units W/m$^2$) at
the receiver location.
\textbf{Significance:} this is Eq.~14 with the $G_\text{Rx}$ and
wavelength-squared factor stripped out --- i.e. just the raw
free-space power density arriving at the receiver's location,
independent of what kind of receiver sits there.

\subsection*{Eq.~(16) --- power via effective aperture}
\begin{equation}
P_\text{Rx} = S_\text{Rx}A_\text{eff}. \tag{16}
\end{equation}
\textbf{Terms:} $A_\text{eff}$ = the Rydberg receiver's effective
aperture (Eq.~17).
\textbf{Significance:} replaces the classical antenna-gain-based
power-capture picture with an aperture-based one --- appropriate since
the Rydberg atoms ``sense'' the field via dipole coupling rather than
guiding currents through a metallic structure.

\subsection*{Eq.~(17) --- effective aperture}
\begin{equation}
A_\text{eff} = \frac{2Z_0N_\text{atoms}\wp_\text{RF}^2\omega_\text{RF}}{\hbar\Gamma_\text{FWHM}}, \qquad \omega_\text{RF}=2\pi f_\text{RF}. \tag{17}
\end{equation}
\textbf{Terms:} $N_\text{atoms}$ = number of atoms effectively
contributing to the reception (\textbf{the paper never gives this a
numeric value or an explicit formula anywhere} --- confirmed by
exhaustive search of the paper and its three cited references; this is
a genuine, real gap, not something you're missing in your reading).
$\Gamma_\text{FWHM}$ = the ``intrinsic EIT linewidth'' (also never given
one single, unambiguous closed-form value that's consistent across
every figure that uses it --- see the per-figure notes near the end of
this document).
\textbf{Significance:} derived in Appendix~B from first principles
(scattering-rate $\times$ energy-per-photon arguments) --- confirms
$A_\text{eff}$ scales with $N_\text{atoms}$ and the transition's own
properties, explicitly \emph{not} the physical cross-sectional area of
the vapor cell. Not used by Fig.~3 (no link budget there).

\subsection*{Eq.~(18) --- LO-free observation model, recast}
\begin{equation}
z = \frac{\wp_\text{RF}}{\hbar}\left|E_\text{RF}hx+n_\text{Ry}\right|. \tag{18}
\end{equation}
\textbf{Terms:} $z$ = the scalar LO-free observation (physically, an
estimated Rabi frequency). $n_\text{Ry}$ = the lumped noise term
covering all relevant noise sources (detailed later in Eq.~46).
\textbf{Significance:} recasts Eq.~13's signal model into the actual
physical quantity the LO-free receiver outputs --- a magnitude
(Rabi-frequency) estimate, not a complex baseband voltage like a
conventional receiver would produce.

\subsection*{Eq.~(19) --- piecewise-constant RF amplitude assumption}
\begin{equation}
x(t) \approx x[k],\quad t\in[kT_m,(k+1)T_m]. \tag{19}
\end{equation}
\textbf{Terms:} $T_m$ = one measurement/scan window duration.
\textbf{Significance:} formalizes the assumption that the incident RF
amplitude doesn't change appreciably during one full coupling-laser
scan --- needed to justify treating one scan as producing one valid
observation.

\subsection*{Eq.~(20) --- coupling-laser scan trajectory}
\begin{equation}
\Delta_c(t) = \Delta_{c,\min}+\frac{\Delta_{c,\max}-\Delta_{c,\min}}{T_m}(t-kT_m),\qquad t\in[kT_m,(k+1)T_m]. \tag{20}
\end{equation}
\textbf{Terms:} $\Delta_{c,\min},\Delta_{c,\max}$ = the scanned detuning
range's endpoints.
\textbf{Significance:} describes the coupling laser being linearly
swept (a ramp) across a fixed detuning window once per measurement
window --- this is exactly the kind of sweep Fig.~3's $\Delta_c$-axis
represents conceptually, though Fig.~3 itself just plots the
steady-state result at each $\Delta_c$, not a literal time-domain scan.

\subsection*{Eq.~(21) --- continuous-time photodetector output}
\begin{equation}
v_k(t) = R_\text{PD}P_\text{out}(\Delta_c(t))+n_{\text{Ry},k}(t),\quad t\in[kT_m,(k+1)T_m]. \tag{21}
\end{equation}
\textbf{Terms:} $R_\text{PD}$ = photodetector responsivity.
$n_{\text{Ry},k}(t)$ = the relevant noise contributions during scan $k$.
\textbf{Significance:} what the PD literally outputs during one scan
--- the real $P_\text{out}(\Delta_c)$ trace (exactly what Fig.~3 plots,
one $\Omega_\text{RF}$ slice at a time) plus noise.

\subsection*{Eq.~(22) --- sampled photodetector output}
\begin{equation}
v_k[m] = R_\text{PD}P_\text{out}(\Delta_c[m])+n_{\text{Ry},k}[m]. \tag{22}
\end{equation}
\textbf{Terms:} $m=0,\dots,M-1$ indexes $M$ discrete detuning sample
points within one scan.
\textbf{Significance:} the discrete-time version of Eq.~21 --- this is
literally what one column of a digitized Fig.~3-style sweep looks like
in practice.

\subsection*{Eq.~(23) --- the LO-free discrete observation}
\begin{equation}
z[k] \triangleq \Omega_\text{RF}[k] = \mathcal E\left(v_k[m]_{m=0}^{M-1}\right). \tag{23}
\end{equation}
\textbf{Terms:} $\mathcal E(\cdot)$ = an AT-splitting extraction
procedure (e.g. peak-separation detection, or model-based fitting).
\textbf{Significance:} formalizes the LO-free receiver's whole job as a
function $\mathcal E$ that turns one entire scan trace (i.e., one
Fig.~3-style curve) into a single scalar $\Omega_\text{RF}$ estimate.
This is also why the LO-free symbol rate is capped at $R_s\le1/T_m$
--- fundamentally, it needs one whole scan per observation.

\section{Section III-B: LO-dressed Rydberg atomic receiver model}

\subsection*{Eq.~(24) --- LO and RF field definitions}
\begin{equation}
E_\text{LO} = A_\text{LO}\cos(2\pi f_\text{LO}t+\phi_\text{LO}), \qquad
E_\text{RF} = A_\text{RF}\cos(2\pi f_\text{RF}t+\phi_\text{RF}). \tag{24a, 24b}
\end{equation}
\textbf{Terms:} $A_\text{LO},A_\text{RF}$ = amplitudes ($A_{(\cdot)}=|E_{(\cdot)}|$).
$f_\text{LO},f_\text{RF}$ = frequencies. $\phi_\text{LO},\phi_\text{RF}$
= phases. Define $\Delta f=f_\text{LO}-f_\text{RF}$,
$\Delta\phi=\phi_\text{LO}-\phi_\text{RF}$.
\textbf{Significance:} simply defines the two real physical RF tones
(the deliberately-injected LO, and the actual signal to be measured)
that get superposed in the LO-dressed architecture.

\subsection*{Eq.~(25) --- exact field superposition}
\begin{align}
E_\text{total} &=E_\text{LO}+E_\text{RF} \notag\\
&=\sqrt{A_\text{LO}^2+A_\text{RF}^2+2A_\text{LO}A_\text{RF}\cos(2\pi\Delta f\,t+\Delta\phi)}\,\cos(2\pi f_\text{LO}t+\phi_\text{LO}). \tag{25}
\end{align}
\textbf{Significance:} the exact sum of the two cosines, rewritten as a
single carrier at $f_\text{LO}$ with a slowly-time-varying \emph{envelope}
amplitude that depends on $\Delta f$ and $\Delta\phi$. This is a pure
trigonometric identity, not an approximation, and shows the physical
origin of envelope/beat detection in the LO-dressed scheme.

\subsection*{Eq.~(26) --- small-signal (Taylor) approximation}
\begin{equation}
E_\text{total} \approx [A_\text{LO}+A_\text{RF}\cos(2\pi\Delta f\,t+\Delta\phi)]\cos(2\pi f_\text{LO}t+\phi_\text{LO}), \qquad (A_\text{RF}/A_\text{LO}\ll1). \tag{26}
\end{equation}
\textbf{Significance:} a second-order Taylor expansion of Eq.~25 around
$A_\text{RF}/A_\text{LO}=0$, valid whenever the RF signal is much weaker
than the injected LO (the intended operating regime). Physically: the
carrier at $f_\text{LO}$ has its amplitude modulated by a slow cosine at
$\Delta f$ --- this is literally amplitude-modulation-style envelope
detection, and is the direct justification for approximating
$\Omega_\text{total}$'s magnitude the way Fig.~6/7's build does.

\subsection*{Eq.~(27) --- closed-form $\rho_{21}$ at resonance, metastable approximation}
\begin{equation}
\rho_{21}\big|_{(\Delta_p,\Delta_c,\Delta_\text{LO},\gamma_3,\gamma_4)=0}
= \frac{j\gamma_2\Omega_p|\Omega_\text{total}|^2}{\gamma_2^2|\Omega_\text{total}|^2+2\Omega_p^2\left(\Omega_c^2+\Omega_p^2+|\Omega_\text{total}|^2\right)}. \tag{27}
\end{equation}
\textbf{Terms:} evaluated at full resonance
($\Delta_p=\Delta_c=\Delta_\text{LO}=0$) and under the approximation
$\gamma_3=\gamma_4=0$ (justified since states 3, 4 are metastable
Rydberg states with much longer lifetimes than state 2, i.e.
$\gamma_3,\gamma_4\ll\gamma_2$). $\Omega_\text{total}=\Omega_\text{LO}+e^{j(2\pi\Delta f\,t+\Delta\phi)}\Omega_\text{RF}$
(complex phasor, defined just before this equation in the text).
\textbf{Significance:} the LO-dressed analogue of Eq.~8 --- another
closed-form steady state, but valid specifically at exact resonance
with the two shortest-lived-state decay rates dropped. This is the
starting point for the entire LDR/distortion derivation in Appendix~C
(Eq.~75--83) and Sec.~IV-B (Eq.~53--58).

\subsection*{Eq.~(28) --- practical LO-dressed output}
\begin{equation}
P_\text{out}(t) = \bar P_0 + \kappa\Omega_\text{RF}\cos(2\pi\Delta f\,t+\Delta\phi). \tag{28}
\end{equation}
\textbf{Terms:} $\bar P_0=P_\text{in}e^{-\alpha}e^{\alpha\Lambda(\Omega_\text{LO},\Gamma)}$
= the average (DC) probe transmission, set purely by the LO bias
$\Omega_\text{LO}$. $\alpha=k_pLC_0\bar A$ = the absorption coefficient
of the probe laser, with $\bar A=\gamma_2\Omega_p/(\gamma_2^2+2\Omega_p^2)$
= the amplitude of the (simplified, three-level) EIT spectrum.
$\Lambda(a,b)=b^2/(b^2+a^2)$ = a normalized Lorentzian function of
variable $a$ with HWHM $b$. $\Gamma=\Omega_p\sqrt{2(\Omega_c^2+\Omega_p^2)/(2\Omega_p^2+\gamma_2^2)}$
= the HWHM of a four-level Lorentzian (defined again, identically, at
Eq.~78). $\kappa=\alpha\bar P_0\kappa_\rho$ = the \textbf{intrinsic gain
coefficient}, with $\kappa_\rho=\partial\Lambda(\Omega_\text{LO},\Gamma)/\partial\Omega_\text{LO}$.
\textbf{Significance:} this is the central result of the LO-dressed
model --- it says $P_\text{out}(t)$ oscillates \emph{linearly} in
$\Omega_\text{RF}$ around a DC bias set entirely by $\Omega_\text{LO}$.
$\kappa$ is the single most important LO-dressed design figure of
merit: the whole point of choosing a good $\Omega_\text{LO}$ is to
maximize $|\kappa|$ (done at Eq.~79). \textbf{Not used by Fig.~3}
(Fig.~3 is LO-free, no $\Omega_\text{LO}$ involved at all).

\subsection*{Eq.~(29) --- AC-coupled readout}
\begin{equation}
P_\text{readout}(t)=P_\text{out}(t)-\bar P_0 = |P(\Delta f)|\cos(2\pi\Delta f\,t+\Delta\phi). \tag{29}
\end{equation}
\textbf{Terms:} $|P(\Delta f)|$ = the amplitude of the single-sided
Fourier spectrum of $P_\text{out}(t)$ evaluated at $\Delta f$ --- i.e.
literally what a spectrum analyzer or lock-in amplifier reads out.
\textbf{Significance:} subtracting off the DC bias $\bar P_0$ isolates
the useful oscillating signal, whose amplitude directly encodes
$\Omega_\text{RF}$ via Eq.~30.

\subsection*{Eq.~(30) --- RF Rabi frequency from the Fourier amplitude}
\begin{equation}
\Omega_\text{RF} = \frac{|P(\Delta f)|}{|\kappa|}. \tag{30}
\end{equation}
\textbf{Significance:} the actual measurement recipe for the
LO-dressed receiver --- measure the Fourier-spectrum peak at $\Delta f$,
divide by the known/calibrated $\kappa$, and you have $\Omega_\text{RF}$
(and hence, via Eq.~12, the RF field amplitude $A_\text{RF}=\hbar\Omega_\text{RF}/\wp_\text{RF}$).
Both $|P(\Delta f)|$ and $|\kappa|$ can be obtained directly from the
optical spectrum, which is why this technique reduces electric-field
measurement to a purely optical-frequency readout.

\section{Section III-C: Noise modeling}

Fig.~2 provides an explicit block-diagram (Tx $\to$ RF signal $\to$
vapor cell $\to$ PD $\to$ LNA $\to$ LIA $\to$ observation) for all three
receivers side by side, labeling exactly where each noise term enters
and with what multiplicative gain. The equations below are the formulas
behind each labeled block.

\subsection*{Eq.~(31) --- AC photocurrent}
\begin{equation}
I_\text{AC} = D\left(|\kappa|\frac{\wp_\text{RF}}{\hbar}\sqrt{P_\text{Rx}}+D|\kappa|\frac{\wp_\text{RF}}{\hbar}n_\text{ex}+Dn_\text{PSN}\right). \tag{31}
\end{equation}
\textbf{Terms:} $D$ = photodiode responsivity (A/W). $n_\text{PSN}$ =
photon shot noise contribution.
\textbf{Significance:} converts the optical AC signal (amplitude
$|\kappa|\Omega_\text{RF}$, per Eq.~28) into an electrical photocurrent,
carrying both the useful signal and the extrinsic + shot-noise
contributions through the same responsivity $D$.

\subsection*{Eq.~(32) --- electrical AC power}
\begin{equation}
P_\text{eAC} = I_\text{AC}^2R_L. \tag{32}
\end{equation}
\textbf{Terms:} $R_L$ = PD output load impedance.
\textbf{Significance:} simple $I^2R$ conversion from photocurrent to
electrical power, at the point right after the photodetector.

\subsection*{Eq.~(33) --- LNA output power}
\begin{equation}
P_\text{LNA} = G_\text{LNA}P_\text{eAC}+\sigma_\text{TN}^2. \tag{33}
\end{equation}
\textbf{Terms:} $G_\text{LNA}$ = low-noise-amplifier gain.
$\sigma_\text{TN}^2$ = thermal (Johnson--Nyquist) noise power, added at
this stage.
\textbf{Significance:} amplifies the PD's electrical output and adds
the amplifier's own thermal noise floor.

\subsection*{Eq.~(34) --- full LO-dressed observation, with channel}
\begin{equation}
z^\text{LO} = \sqrt{P_\text{Rx}G_\text{LNA}R_L}\,D|\kappa|\frac{\wp_\text{RF}}{\hbar}hx+n_{\text{Ry,LO}}. \tag{34}
\end{equation}
\textbf{Significance:} the complete end-to-end LO-dressed observation
model, folding the wireless channel $h$, the transmitted symbol $x$, and
every gain stage (PD responsivity, LNA gain, load resistance) into one
expression, with $n_{\text{Ry,LO}}$ collecting all noise (detailed at
Eq.~47). This is the LO-dressed counterpart of Eq.~18.

\subsection*{Eq.~(35) --- noise-equivalent field (NEF)}
\begin{equation}
\text{NEF}=\frac{|E_\text{RF}|_\text{min}}{\sqrt{B}}. \tag{35}
\end{equation}
\textbf{Terms:} $|E_\text{RF}|_\text{min}$ = minimum detectable field
amplitude. $B$ = measurement bandwidth.
\textbf{Significance:} a bandwidth-normalized noise floor expressed in
the receiver's natural ``field'' units (V/m/$\sqrt{\text{Hz}}$) --- a
smaller NEF means a more sensitive receiver. This is the standard
figure of merit quoted throughout the Rydberg-sensor literature.

\subsection*{Eq.~(36) --- extrinsic NEF}
\begin{equation}
\text{NEF}_\text{ex} = \sqrt{\frac{16\pi f_\text{RF}^2}{3\epsilon_0c^3}\Theta(f_\text{RF},T)}. \tag{36}
\end{equation}
\textbf{Terms:} $c$ = speed of light.
\textbf{Significance:} the external (BBR + vacuum-fluctuation) noise
floor, expressed as an NEF --- this is a fundamental limit that no
free-space receiver, Rydberg or classical, can beat.

\subsection*{Eq.~(37) --- modified Callen--Welton spectral factor}
\begin{equation}
\Theta(f_\text{RF},T) = \hbar f_\text{RF}
\begin{cases}
\tfrac12 n_\text{th}(f_\text{RF},T)+\tfrac12, & \text{homodyne},\\
2n_\text{th}(f_\text{RF},T)+1, & \text{heterodyne}.
\end{cases} \tag{37}
\end{equation}
\textbf{Significance:} distinguishes the noise contribution depending
on whether the detection is homodyne or heterodyne --- this factor
folds both the thermal (BBR) and quantum vacuum contributions into one
expression via $n_\text{th}$ (Eq.~38).

\subsection*{Eq.~(38) --- thermal photon occupation number}
\begin{equation}
n_\text{th}(f_\text{RF},T) = \frac{1}{\exp(\hbar f_\text{RF}/k_BT)-1}. \tag{38}
\end{equation}
\textbf{Significance:} the standard Bose--Einstein thermal occupation
number at frequency $f_\text{RF}$ and temperature $T$ --- the number of
thermal photons per mode.

\subsection*{Eq.~(39) --- extrinsic noise power}
\begin{equation}
\sigma_\text{ex}^2 = \frac{\text{NEF}_\text{ex}^2}{2Z_0}A_\text{eff}B. \tag{39}
\end{equation}
\textbf{Significance:} converts the extrinsic NEF (a field-domain
quantity) into an actual noise \emph{power} at the receiver, using the
effective aperture $A_\text{eff}$ (Eq.~17) to do the field$\to$power
conversion and $B$ to set the bandwidth.

\subsection*{Eq.~(40) --- standard quantum limit (SQL) field}
\begin{equation}
\frac{E_\text{SQL}}{\sqrt{\text{Hz}}} = \frac{\hbar}{\wp_\text{RF}\sqrt{N_\text{atoms}T_2}}. \tag{40}
\end{equation}
\textbf{Terms:} $T_2$ = the EIT coherence time.
\textbf{Significance:} the fundamental atomic-shot-noise-limited
sensitivity floor --- scales down as $1/\sqrt{N_\text{atoms}}$, the
standard quantum-limit scaling for $N_\text{atoms}$ independent quantum
probes.

\subsection*{Eq.~(41) --- quantum-projection-noise (QPN) power}
\begin{equation}
\sigma_\text{QPN}^2 = \left(\frac{E_\text{SQL}}{\sqrt{\text{Hz}}}\right)^2\frac{A_\text{eff}B}{2Z_0}. \tag{41}
\end{equation}
\textbf{Significance:} same field-to-power conversion pattern as
Eq.~39, but starting from the SQL field instead of the extrinsic NEF ---
this is the fundamental atomic-physics noise contribution (as opposed
to the BBR-driven extrinsic noise).

\subsection*{Eq.~(42) --- average photocurrent from probe power}
\begin{equation}
\bar I = \frac{\eta_\text{eff}e}{2\pi\hbar f_p}|P_\text{out}|. \tag{42}
\end{equation}
\textbf{Terms:} $\eta_\text{eff}\in[0,1]$ = PD quantum efficiency.
$e$ = elementary charge. $f_p$ = probe laser frequency.
\textbf{Significance:} standard photodiode relation --- converts
incident optical power $|P_\text{out}|$ into average photocurrent via
the quantum efficiency and photon energy $\hbar f_p$.

\subsection*{Eq.~(43) --- photon shot noise (PSN) power}
\begin{equation}
\sigma_\text{PSN}^2 = 2eB\bar I. \tag{43}
\end{equation}
\textbf{Significance:} the standard shot-noise formula for a
photocurrent $\bar I$ over bandwidth $B$ --- an intrinsic, unavoidable
consequence of the discreteness of photon detection.

\subsection*{Eq.~(44) --- thermal (Johnson--Nyquist) noise power}
\begin{equation}
\sigma_\text{TN}^2 =
\begin{cases}
4k_BTB, & \text{TIA front-end},\\
Fk_BTB, & \text{LNA front-end}.
\end{cases} \tag{44}
\end{equation}
\textbf{Terms:} $F$ = LNA noise factor.
\textbf{Significance:} the electronic thermal noise floor, with two
different forms depending on whether the front end is a
transimpedance amplifier (dominated by feedback-resistor Johnson noise)
or a matched LNA (dominated by the noise factor $F$).

\subsection*{Eq.~(45) --- LO-free small-signal transduction gain}
\begin{equation}
G_\text{Ry} = \frac{\partial P_\text{out}}{\partial E_\text{RF}}. \tag{45}
\end{equation}
\textbf{Significance:} defines the LO-free receiver's own intrinsic
gain --- the local slope of the probe transmission with respect to the
incident RF field. This is deliberately analogous to $\kappa$
(Eq.~28) for the LO-dressed case, letting the two architectures be
compared apples-to-apples in the noise formulas below.

\subsection*{Eq.~(46) --- total LO-free noise power}
\begin{equation}
\sigma_\text{Ry}^2 = G_\text{Ry}^2D^2G_\text{LNA}^2\sigma_\text{ex}^2+D^2G_\text{LNA}^2\sigma_\text{QPN}^2+G_\text{LNA}^2\sigma_\text{PSN}^2+\sigma_\text{TN}^2. \tag{46}
\end{equation}
\textbf{Significance:} sums all four noise contributions (extrinsic,
QPN, PSN, TN), each referred forward through the correct chain of gains
up to the final observation point. Note the extrinsic term is scaled
by $G_\text{Ry}^2$, which the text calls ``modest,'' since the LO-free
architecture ``operates on the natural slope of an EIT doublet.''

\subsection*{Eq.~(47) --- total LO-dressed noise power}
\begin{equation}
\sigma_\text{Ry,LO}^2 = \kappa^2D^2G_\text{LNA}^2\sigma_\text{ex}^2+D^2G_\text{LNA}^2\sigma_\text{QPN}^2+G_\text{LNA}^2\sigma_\text{PSN}^2+\sigma_\text{TN}^2. \tag{47}
\end{equation}
\textbf{Significance:} \textbf{the single most important
noise-comparison equation in the paper} --- structurally identical to
Eq.~46, but the extrinsic term is scaled by $\kappa^2$ instead of
$G_\text{Ry}^2$. Since $\kappa$ can be one to two orders of magnitude
larger than $G_\text{Ry}$ (stated directly in the text right after
Eq.~47), this single substitution is the quantitative origin of the
LO-dressed receiver's SNR advantage seen later in Fig.~8/9/10.

\section{Section IV-A: SNR performance}

\subsection*{Eq.~(48) --- the dimensionless resolvability ratio}
\begin{equation}
\mathcal R \equiv \frac{\Omega_\text{RF}}{\Gamma_\text{FWHM}}. \tag{48}
\end{equation}
\textbf{Significance:} the single dimensionless figure of merit for
LO-free AT-splitting resolvability --- how large the RF coupling is
relative to the intrinsic linewidth. \textbf{This is exactly Fig.~4's
$y$-axis.} Achieving $\mathcal R\gg1$ is stated as a necessary condition
for clear splitting and high-fidelity LO-free sensing.

\subsection*{Eq.~(49) --- Fisher-information penalty}
\begin{equation}
G(\mathcal R) = \frac{\mathcal R^2}{1+\frac{1}{2\mathcal R^2}}. \tag{49}
\end{equation}
\textbf{Significance:} derived rigorously in Appendix~D from Poisson
statistics and the Cram\'er--Rao bound (not an ad-hoc fudge factor) ---
quantifies the information penalty/resolvability loss of trying to
estimate $\Omega_\text{RF}$ from an AT-split spectrum. Collapses toward
0 as $\mathcal R\to0$ (matching the qualitative statement that below
$\mathcal R=1$ the splitting is unobservable), approaches $\mathcal R^2$
for $\mathcal R\gg1$.

\subsection*{Eq.~(50) --- LO-free SNR}
\begin{equation}
\text{SNR}_\text{Ry} = \frac{P_\text{Rx}|h|^2}{\sigma_\text{Ry}^2}\left(\frac{\wp_\text{RF}}{\hbar\Gamma_\text{FWHM}}\right)^2G(\mathcal R). \tag{50}
\end{equation}
\textbf{Significance:} integrates the classical wireless link budget
($P_\text{Rx}|h|^2/\sigma_\text{Ry}^2$) with the RF-to-atomic conversion
gain $(\wp_\text{RF}/\hbar\Gamma_\text{FWHM})^2$ and the resolvability
penalty $G(\mathcal R)$ --- these three multiplicative factors are
individually dimensionless-consistent, so $G(\mathcal R)$ only reshapes
the numerical value, not the physical units.

\subsection*{Eq.~(51) --- LO-free sensitivity floor}
\begin{equation}
E_\text{min}^\text{strong} = \frac{\hbar\Gamma_\text{EIT}}{\wp_\text{RF}}\left(\frac{2Z_0\sigma_\text{Ry}^2}{A_\text{eff}|h|^2}\right)^{1/4}. \tag{51}
\end{equation}
\textbf{Significance:} obtained by setting $\text{SNR}_\text{Ry}=1$ in
Eq.~50 and solving for the minimum detectable field --- the paper
reports a computed value of $38.15\,\mu\text{V/cm}/\sqrt{\text{Hz}}$,
cross-checked as consistent with existing literature.

\subsection*{Eq.~(52) --- LO-dressed SNR}
\begin{equation}
\text{SNR}_\text{Ry,LO} = \frac{G_\text{LNA}R_LD^2\kappa^2\frac{\wp_\text{RF}^2}{\hbar^2}P_\text{Rx}|h|^2}{\sigma_\text{Ry,LO}^2}. \tag{52}
\end{equation}
\textbf{Significance:} the LO-dressed SNR expression, built from the
signal model (Eq.~34) and total noise (Eq.~47) --- note $\kappa^2$
appears explicitly in the numerator (signal power) as well as inside
$\sigma_\text{Ry,LO}^2$'s extrinsic term (Eq.~47), which is why Eq.~61
below is a more physically revealing rearrangement.

\section{Section IV-B: Distortion effect / linear dynamic range}

\subsection*{Eq.~(53) --- restructured probe transmission}
\begin{equation}
P_\text{out} = P_\text{in}e^{-\alpha}e^{\alpha\Lambda(\Omega_\text{total},\Lambda)}. \tag{53}
\end{equation}
\textbf{Significance:} rewrites $P_\text{out}$ (same physical quantity
as Eq.~2/28) explicitly as a function of the Lorentzian $\Lambda$
evaluated at the \emph{total} Rabi frequency $\Omega_\text{total}$ ---
the starting point for the Taylor-expansion-based distortion analysis
that follows.

\subsection*{Eq.~(54) --- Taylor expansion of the Lorentzian}
\begin{equation}
\Lambda(\Omega_\text{total},\Lambda) = \Lambda_0+\Lambda_0'\delta\Omega+\frac12\Lambda_0''\delta\Omega^2+\frac16\Lambda_0'''\delta\Omega^3+\dots \tag{54}
\end{equation}
\textbf{Terms:} $\delta\Omega=\Omega_\text{total}-\Omega_\text{LO}$.
$\Lambda_0'=\dfrac{-2\Gamma^2\Omega_\text{LO}}{(\Gamma^2+\Omega_\text{LO}^2)^2}$,
$\Lambda_0''=\dfrac{-2\Gamma^2(\Gamma^2-3\Omega_\text{LO}^2)}{(\Gamma^2+\Omega_\text{LO}^2)^3}$
(explicit derivative formulas, evaluated at the bias point
$\Omega_\text{LO}$).
\textbf{Significance:} a standard Taylor expansion around the LO bias
point --- this is exactly the mathematical machinery used to separate
$P_\text{out}$ into a desired linear term plus 2nd- and 3rd-order
distortion terms in Eq.~55.

\subsection*{Eq.~(55) --- $P_\text{out}$ split into desired + distortion terms}
\begin{equation}
P_\text{out} =\bar P_0 + \underbrace{\kappa\Omega_\text{RF}\cos\theta}_{\text{desired}}
+ \underbrace{\tfrac14\alpha\bar P_0\Lambda_0''\Omega_\text{RF}^2(1+\cos2\theta)}_{\text{2nd-order distortion}}
+ \underbrace{\tfrac16\alpha\bar P_0\Lambda_0'''\Omega_\text{RF}^3\cos^3\theta}_{\text{3rd-order distortion}}. \tag{55}
\end{equation}
\textbf{Terms:} $\theta=2\pi\Delta f\,t+\Delta\phi$.
\textbf{Significance:} substituting Eq.~54 into $P_\text{out}$ and
retaining terms up to third order in $\Omega_\text{RF}$ produces exactly
three physically distinct terms: the wanted linear signal (oscillating
at $\Delta f$, amplitude $\propto\kappa$), a second-harmonic distortion
term (oscillating at $2\Delta f$, amplitude $\propto\Lambda_0''$), and a
third-harmonic distortion term (oscillating at $3\Delta f$, amplitude
$\propto\Lambda_0'''$). This is the direct mathematical origin of the
``linear dynamic range'' concept: as $\Omega_\text{RF}$ grows, the
distortion terms (which grow as $\Omega_\text{RF}^2,\Omega_\text{RF}^3$)
eventually catch up to and overwhelm the desired term (which only grows
as $\Omega_\text{RF}^1$).

\subsection*{Eq.~(56) --- total-harmonic-distortion (THD) ratio}
\begin{equation}
\frac{P(\Delta f)}{P(2\Delta f)} = \frac{\kappa\Omega_\text{RF}}{\tfrac14\alpha\bar P_0\Lambda_0''\Omega_\text{RF}^2} \le \epsilon. \tag{56}
\end{equation}
\textbf{Terms:} $\epsilon$ = the total harmonic distortion (THD)
tolerance. \textbf{Note:} the paper's own equation number/label here
reads ``$P(\Delta f)/P(3\Delta f)$'' in one place, which appears to be a
typo, since the actual formula shown uses the second-harmonic term
$P(2\Delta f)=\frac14\alpha\bar P_0\Lambda_0''\Omega_\text{RF}^2$
defined in the sentence immediately preceding it --- treat this as
comparing fundamental to \emph{second}-harmonic power, consistent with
Eq.~55's ``2nd-order distortion'' label.
\textbf{Significance:} defines the design criterion --- keep the
second-harmonic-to-fundamental power ratio below some tolerance
$\epsilon$ --- that Eq.~57 directly solves for $\Omega_\text{RF}$.

\subsection*{Eq.~(57) --- second-order LDR upper bound (VERIFY: contains $\epsilon$)}
\begin{equation}
\Omega_\text{RF,max}'' = \frac{4\epsilon|\Lambda_0'|}{|\Lambda_0''|}. \tag{57}
\end{equation}
\textbf{Terms:} $\epsilon$ = the same THD tolerance from Eq.~56 ---
\textbf{confirmed by directly re-reading the rendered paper page: this
equation DOES contain $\epsilon$.} It is not simply
$4|\Lambda_0'/\Lambda_0''|$ --- that ratio is only the coefficient that
$\epsilon$ multiplies.
\textbf{Significance:} solving Eq.~56 for $\Omega_\text{RF}$ gives the
maximum RF Rabi frequency you can apply before second-order distortion
exceeds your chosen tolerance $\epsilon$. \textbf{Practical caveat}:
$\Lambda_0''\to0$ exactly at the true optimal bias point
$\Omega_\text{LO,opt}$ (Eq.~79), which makes this bound blow up
(divide by zero) there --- this is not a flaw in the equation, it's the
direct mathematical statement that ``second-order distortion vanishes
at the optimum,'' which is precisely why Eq.~58 (third-order) is needed
at that specific bias point.

\subsection*{Eq.~(58) --- third-order LDR upper bound (used at the true optimum)}
\begin{equation}
\Omega_\text{RF,max}''' = \sqrt{\frac{6\epsilon|\Lambda_0'|}{|\Lambda_0'''|}}. \tag{58}
\end{equation}
\textbf{Terms:} $\Lambda_0'''$ = the third derivative of $\Lambda$ with
respect to $\Omega_\text{LO}$ at the bias point (not given an explicit
closed form in the paper text shown, but obtainable by differentiating
$\Lambda_0''$'s formula once more, or numerically from real data).
\textbf{Significance:} the governing LDR bound specifically \emph{at}
the optimal bias $\Omega_\text{LO,opt}$, where second-order distortion
has vanished and third-order distortion becomes the limiting factor.
\textbf{The paper never assigns $\epsilon$ a numeric value anywhere} ---
it calls $\epsilon$ ``the designer-specified tolerance,'' explicitly
framing it as a free engineering choice, not a physical constant. Any
numeric LDR width computed from Eq.~57 or Eq.~58 therefore requires
choosing your own $\epsilon$ and labeling it as such.

\section{Section IV-C / V-B: mutual information}

\subsection*{Eq.~(59) --- LO-free mutual information}
\begin{equation}
\mathcal I(z;x) = \frac12\ln\left(\frac{\pi(\mathfrak r+1)}{2}\right)-\ln I_0(\mathfrak r)+\mathfrak r\frac{I_1(\mathfrak r)}{I_0(\mathfrak r)}. \tag{59}
\end{equation}
\textbf{Terms:} $I_0,I_1$ = modified Bessel functions of the first
kind, orders 0 and 1. $\mathfrak r$ = the Rician factor, approximated by
$\mathfrak r\approx\text{SNR}_\text{Ry}$.
\textbf{Significance:} the theoretical mutual information for the
LO-free receiver's signal model (Eq.~18), which behaves like a
Rician-distributed amplitude observation --- this is exactly what
Fig.~9's green ``LO-free'' curve plots.

\subsection*{Eq.~(60) --- LO-dressed mutual information}
\begin{equation}
\mathcal I(z^\text{LO};x) = e^{1/\text{SNR}_\text{Ry,LO}}E_1(1/\text{SNR}_\text{Ry,LO})\ln2. \tag{60}
\end{equation}
\textbf{Terms:} $E_1(x)=\int_x^\infty\frac{e^{-t}}{t}dt$ = the
exponential integral function.
\textbf{Significance:} the LO-dressed mutual information, based on the
signal model in Eq.~34 (a complex Gaussian observation, unlike the
LO-free receiver's real-amplitude-only observation) --- this is what
Fig.~9's orange ``LO-dressed'' curves plot.

\subsection*{Eq.~(61) --- restructured LO-dressed SNR}
\begin{equation}
\text{SNR}_\text{Ry,LO} = \cfrac{R_L\frac{\wp_\text{RF}^2}{\hbar^2}P_\text{Rx}|h|^2}{G_\text{LNA}\sigma_\text{ex}^2+\frac{G_\text{LNA}\sigma_\text{QPN}^2}{\kappa^2}+\frac{G_\text{LNA}\sigma_\text{PSN}^2}{D^2\kappa^2}+\frac{\sigma_\text{TN}^2}{D^2\kappa^2G_\text{LNA}}}. \tag{61}
\end{equation}
\textbf{Significance:} algebraically identical to Eq.~52, but
rearranged to make the mechanism explicit: the common extrinsic
(BBR) noise $\sigma_\text{ex}^2$ is \emph{not} divided by $\kappa^2$
(same $\kappa$-dependence as the signal itself, so it doesn't help),
while QPN, PSN, and TN all \emph{are} divided by $\kappa^2$ ---
explaining precisely why the LO-dressed advantage is concentrated in
non-BBR-limited operating regimes, not a violation of any shared noise
floor.

\section{Appendix A: solving the master equation}

\subsection*{Eq.~(62) --- the $\rho_{21}$ commutator element}
\begin{equation}
[\mathbf H,\boldsymbol\rho]_{21} = \frac{\hbar\Omega_p}{2}(\rho_{11}-\rho_{22})-\hbar\Delta_p\rho_{21}+\frac{\hbar\Omega_c}{2}\rho_{31}. \tag{62}
\end{equation}
\textbf{Significance:} the explicit $(2,1)$ matrix element of the
commutator $[\mathbf H,\boldsymbol\rho]$, computed directly from
Eq.~6's Hamiltonian using the general rule
$[\mathbf H,\boldsymbol\rho]_{ij}=\sum_k H_{ik}\rho_{kj}-\rho_{ik}H_{kj}$
--- this is the first concrete step toward deriving Eq.~8.

\subsection*{Eq.~(63) --- $\rho_{21}$ evolution equation}
\begin{equation}
\dot\rho_{21} = -j\frac{\Omega_p}{2}(\rho_{11}-\rho_{22})+j\Delta_p\rho_{21}-j\frac{\Omega_c}{2}\rho_{31}-\gamma_{21}\rho_{21}. \tag{63}
\end{equation}
\textbf{Significance:} substitutes Eq.~62 into the master equation
(Eq.~5) and adds the dephasing term $-\gamma_{21}\rho_{21}$ from the
Lindblad part (Eq.~7) --- the first of three coupled differential
equations (with Eq.~64a/b) needed to solve for the full steady state.

\subsection*{Eq.~(64) --- $\rho_{31}$ and $\rho_{41}$ evolution equations}
\begin{align}
\dot\rho_{31} &= -j\frac{\Omega_c}{2}(\rho_{21}-\rho_{32})+j(\Delta_p+\Delta_c)\rho_{31}-j\frac{\Omega_\text{RF}}{2}\rho_{41}-\gamma_{31}\rho_{31}, \tag{64a}\\
\dot\rho_{41} &= -j\frac{\Omega_\text{RF}}{2}(\rho_{31}-\rho_{42})+j(\Delta_p+\Delta_c+\Delta_\text{RF})\rho_{41}-\gamma_{41}\rho_{41}. \tag{64b}
\end{align}
\textbf{Significance:} the analogous evolution equations for the two
other relevant coherences $\rho_{31},\rho_{41}$, obtained the same way
as Eq.~63 --- together, Eq.~63/64a/64b form the full coupled system that
needs to be solved (at steady state) to get $\rho_{21}$.

\subsection*{Eq.~(65) --- steady-state equations after approximations}
\begin{align}
(\gamma_{21}-j\Delta_p)\rho_{21} &= -j\frac{\Omega_p}{2}(\rho_{11}-\rho_{22})-j\frac{\Omega_c}{2}\rho_{31}, \tag{65a}\\
[\gamma_{31}-j(\Delta_p+\Delta_c)]\rho_{31} &= -j\frac{\Omega_c}{2}\rho_{21}-j\frac{\Omega_\text{RF}}{2}\rho_{41}, \tag{65b}\\
[\gamma_{41}-j(\Delta_p+\Delta_c+\Delta_\text{RF})]\rho_{41} &= -j\frac{\Omega_\text{RF}}{2}\rho_{31}. \tag{65c}
\end{align}
\textbf{Significance:} sets $\dot\rho_{21}=\dot\rho_{31}=\dot\rho_{41}=0$
(the steady-state condition) in Eq.~63/64a/64b, and applies two stated
approximations: (i) $\rho_{11}\approx1$ (ground-state population
probability), so $\rho_{22},\rho_{33},\rho_{44}$ are small and neglected;
(ii) $\rho_{32},\rho_{42}\to0$ (their influence on $\rho_{21}$ is
assumed marginal). This produces a clean, linear, coupled 3-equation
system.

\subsection*{Eq.~(66) --- $\rho_{31}$ eliminated}
\begin{equation}
\left[\gamma_{31}-j(\Delta_p+\Delta_c)-\frac{(\Omega_\text{RF}/2)^2}{\gamma_{41}-j(\Delta_p+\Delta_c+\Delta_\text{RF})}\right]\rho_{31} = -j\frac{\Omega_c}{2}\rho_{21}. \tag{66}
\end{equation}
\textbf{Significance:} substitutes Eq.~65c into Eq.~65b to eliminate
$\rho_{41}$, leaving a single equation relating $\rho_{31}$ to
$\rho_{21}$ --- the ``innermost'' layer of the continued fraction that
will appear in Eq.~8.

\subsection*{Eq.~(67) --- $\rho_{21}$ fully resolved (immediately gives Eq.~8)}
\begin{equation}
(\gamma_{21}-j\Delta_p)\rho_{21} = -j\frac{\Omega_p}{2}(\rho_{11}-\rho_{22})-j\frac{\Omega_c}{2}\left(\cfrac{-j\frac{\Omega_c}{2}}{\gamma_{31}-j(\Delta_p+\Delta_c)-\cfrac{(\Omega_\text{RF}/2)^2}{\gamma_{41}-j(\Delta_p+\Delta_c+\Delta_\text{RF})}}\rho_{21}\right). \tag{67}
\end{equation}
\textbf{Significance:} substitutes Eq.~66 into Eq.~65a. Letting
$\rho_{11}\to1,\rho_{22}\to0$ (weak-probe limit) and solving for
$\rho_{21}$ produces exactly Eq.~8 --- this equation is the last
algebraic step of that derivation, shown explicitly so the full chain
from the Hamiltonian to the closed form is traceable.

\section{Appendix B: deriving the effective aperture $A_\text{eff}$}

\subsection*{Eq.~(68) --- steady-state excited-state population}
\begin{equation}
\rho_{ee} = \frac{(\Omega_\text{RF}/\Gamma_\text{FWHM})^2}{1+4(\Delta_c/\Gamma_\text{FWHM})^2+2(\Omega_\text{RF}/\Gamma_\text{FWHM})^2}. \tag{68}
\end{equation}
\textbf{Terms:} $\rho_{ee}$ = the excited-state (Rydberg) population
fraction, cited from a standard two-level-atom saturation formula
[Steck, Rb87 D-line data].
\textbf{Significance:} models how strongly the atoms are driven by the
RF field, in a simplified two-level picture, as a function of the ratio
$\Omega_\text{RF}/\Gamma_\text{FWHM}$ and detuning $\Delta_c/\Gamma_\text{FWHM}$
--- the starting point for computing a scattering rate.

\subsection*{Eq.~(69) --- total scattering rate}
\begin{equation}
\bar R = \Gamma_\text{FWHM}\rho_{ee}. \tag{69}
\end{equation}
\textbf{Significance:} the rate at which the atom absorbs and
re-emits RF photons --- population fraction times the natural
linewidth (standard two-level scattering-rate relation).

\subsection*{Eq.~(70) --- near-resonant, weak-drive scaling}
\begin{equation}
\bar R = \frac{\Omega_\text{RF}^2}{\Gamma_\text{FWHM}} = \frac{\wp_\text{RF}^2E_\text{RF}^2}{\hbar^2\Gamma_\text{FWHM}}, \qquad (\Delta_c\approx0,\ \Omega_\text{RF}\ll\Gamma_\text{FWHM}). \tag{70}
\end{equation}
\textbf{Significance:} simplifies Eq.~68/69 in the specific limit
relevant to weak RF fields at line center, and substitutes
$\Omega_\text{RF}=\wp_\text{RF}E_\text{RF}/\hbar$ (Eq.~12) to express the
scattering rate directly in terms of the field.

\subsection*{Eq.~(71) --- power absorbed per atom}
\begin{equation}
P_\text{atom} = \hbar\omega_\text{RF}\bar R = \frac{\wp_\text{RF}^2E_\text{RF}^2\omega_\text{RF}}{\hbar\Gamma_\text{FWHM}}. \tag{71}
\end{equation}
\textbf{Significance:} each absorption--re-emission event removes one
photon's worth of energy $\hbar\omega_\text{RF}$ from the field ---
multiplying by the scattering rate $\bar R$ gives the power drawn from
the field by a single atom.

\subsection*{Eq.~(72) --- total power absorbed by all atoms}
\begin{equation}
P_\text{total} = N_\text{atoms}P_\text{atom} = \frac{N_\text{atoms}\wp_\text{RF}^2E_\text{RF}^2\omega_\text{RF}}{\hbar\Gamma_\text{FWHM}}. \tag{72}
\end{equation}
\textbf{Significance:} simply scales the single-atom result by
$N_\text{atoms}$, assuming independent, non-interacting atoms (already
justified via the blockade/Rydberg-fraction argument discussed under
Fig.~4/Sec.~V-A).

\subsection*{Eq.~(73) --- matching to a flux $\times$ aperture form}
\begin{equation}
\frac{E_\text{RF}^2}{2Z_0}A_\text{eff} = \frac{N_\text{atoms}\wp_\text{RF}^2E_\text{RF}^2\omega_\text{RF}}{\hbar\Gamma_\text{FWHM}}. \tag{73}
\end{equation}
\textbf{Significance:} defines the effective aperture $A_\text{eff}$
implicitly, by demanding that (power flux density) $\times$
$A_\text{eff}$ reproduce the total absorbed power computed in Eq.~72 ---
the standard way an ``aperture'' is defined for any receiver.

\subsection*{Eq.~(74) --- effective aperture, final form}
\begin{equation}
A_\text{eff} = \frac{2Z_0N_\text{atoms}\wp_\text{RF}^2\omega_\text{RF}}{\hbar\Gamma_\text{FWHM}}. \tag{74}
\end{equation}
\textbf{Significance:} solving Eq.~73 for $A_\text{eff}$ reproduces
Eq.~17 exactly --- confirms the two are the same formula, one derived
from first principles here, one just stated earlier in the main text.

\section{Appendix C: deriving the LO-dressed optimum}

\subsection*{Eq.~(75) --- restructured $\Im[\rho_{21}]$}
\begin{equation}
\Im[\rho_{21}(\Omega_\text{RF})] = \bar A[1-\Lambda(\Omega_\text{LO},\Gamma)-\kappa_\rho\Omega_\text{RF}\cos(2\pi\Delta f\,t+\Delta\phi)]. \tag{75}
\end{equation}
\textbf{Significance:} rewrites Eq.~27's $\Im[\rho_{21}]$ using the
Lorentzian $\Lambda$ notation, exposing $\kappa_\rho=\partial\Lambda(\Omega_\text{LO},\Gamma)/\partial\Omega_\text{LO}$
(the intrinsic gain coefficient associated with the coherence itself,
before the Beer--Lambert exponential is applied) as a clean, separate
factor.

\subsection*{Eq.~(76) --- $\Im[\rho_{21}]$ at $\Omega_\text{RF}=0$, versus $\Delta_c$}
\begin{equation}
\Im[\rho_{21}(\Delta_c)]\big|_{(\Delta_p,\Delta_\text{LO},\gamma_3,\gamma_4,\Omega_\text{RF},\Omega_\text{LO})=0}
= \bar A-\bar A\,\Lambda\!\left(\Delta_c,\Gamma_\text{HWHM}^{(3)}\right). \tag{76}
\end{equation}
\textbf{Significance:} the special case with no RF/LO field at all ---
recovers the plain three-level EIT lineshape as a function of $\Delta_c$
alone, expressed via the three-level HWHM $\Gamma_\text{HWHM}^{(3)}$
(Eq.~77). This is a useful cross-check: it is a completely independent
derivation of the basic (no-RF) EIT peak shape, separate from the
4-level Fig.~3 machinery.

\subsection*{Eq.~(77) --- three-level HWHM}
\begin{equation}
\Gamma_\text{HWHM}^{(3)} = \frac{\Omega_c^2+\Omega_p^2}{2\sqrt{\gamma_2^2+2\Omega_p^2}}. \tag{77}
\end{equation}
\textbf{Significance:} the half-width at half-maximum of the
\emph{three}-level (no RF/Rydberg-coupling) EIT spectrum --- purely a
function of the two laser Rabi frequencies and the intermediate-state
decay rate $\gamma_2$. \textbf{Caution, already established in this
project}: this quantity, used alone, is \emph{not} the correct
$\Gamma_\text{FWHM}$ for Fig.~4/Eq.~48's LO-free normalization (that
candidate was tested and found to predict AT splitting at $\mathcal
R\approx0.3$--0.4, well below the paper's own stated $\mathcal R=1$
threshold).

\subsection*{Eq.~(78) --- four-level $\Gamma$}
\begin{equation}
\Gamma = \frac{2\sqrt2\,\Omega_p}{\sqrt{\Omega_c^2+\Omega_p^2}}\,\Gamma_\text{HWHM}^{(3)}. \tag{78}
\end{equation}
\textbf{Significance:} rescales the three-level HWHM into the
\emph{four}-level $\Gamma$ that appears throughout the LO-dressed
Lorentzian $\Lambda(\Omega_\text{LO},\Gamma)$ (Eq.~28, 53, 54). This is
the $\Gamma$ used for Fig.~6/7's $\Omega_\text{LO,opt}$ (Eq.~79) ---
numerically, with the paper's stated $\Omega_p/2\pi=8$~MHz,
$\Omega_c/2\pi=1$~MHz, $\gamma_2/2\pi=5.2$~MHz, this evaluates to
$\Gamma=7.3255$~MHz, giving $\Omega_\text{LO,opt}=4.2294$~MHz ---
matching the paper's stated 4.23~MHz almost exactly. \textbf{This
$\Gamma$ is a structurally different quantity from
$\Gamma_\text{FWHM}$} used in the LO-free Eq.~17/48/50 --- do not
conflate the two across figures.

\subsection*{Eq.~(79) --- the optimal LO bias}
\begin{equation}
\Omega_\text{LO,opt} = \frac{\sqrt3}{3}\Gamma, \qquad (\kappa_\rho)_\text{max}=\frac{3\sqrt3}{8\Gamma}. \tag{79}
\end{equation}
\textbf{Significance:} found by solving
$\partial^2\Lambda(\Omega_\text{LO},\Gamma)/\partial\Omega_\text{LO}^2=0$
--- the single, physically meaningful bias point that simultaneously
maximizes the intrinsic gain coefficient $\kappa_\rho$ and (by
definition of being a stationary point of the second derivative)
eliminates second-order distortion. This is the vertical line plotted
in Fig.~6(a) and the basis for Fig.~7's ``fixed $\kappa$'' curve.

\subsection*{Eq.~(80) --- practical $P_\text{out}(t)$, full derivation chain}
\begin{align}
P_\text{out}(t)
&= P_\text{in}e^{-k_pLC_0\Im(\rho_{21})} \notag\\
&= P_\text{in}e^{-k_pLC_0\bar A[1-\Lambda(\Omega_\text{LO},\Gamma)-\kappa_\rho\Omega_\text{RF}\cos(2\pi\Delta f\,t+\Delta\phi)]} \notag\\
&= P_\text{in}e^{-\alpha}e^{\alpha\Lambda(\Omega_\text{LO},\Gamma)}e^{\alpha\kappa_\rho\Omega_\text{RF}\cos(2\pi\Delta f\,t+\Delta\phi)} \notag\\
&= P_\text{in}e^{-\alpha}e^{\alpha\Lambda(\Omega_\text{LO},\Gamma)}e^{\kappa\Omega_\text{RF}\cos(2\pi\Delta f\,t+\Delta\phi)}. \tag{80}
\end{align}
\textbf{Significance:} shows step-by-step how substituting Eq.~75 into
the Beer--Lambert law (Eq.~2, with $\chi=C_0\rho_{21}$) and factoring
out constants leads to the exponential form used just before the final
linearization (Eq.~81). Each line is an exact algebraic rearrangement,
no approximation yet.

\subsection*{Eq.~(81) --- linearized practical output}
\begin{equation}
P_\text{out}(t) \approx P_\text{in}e^{-\alpha}e^{\alpha\Lambda(\Omega_\text{LO},\Gamma)}(1+\alpha\kappa_\rho\Omega_\text{RF}\cos(2\pi\Delta f\,t+\Delta\phi)) = \bar P_0+\kappa\Omega_\text{RF}\cos(2\pi\Delta f\,t+\Delta\phi). \tag{81}
\end{equation}
\textbf{Significance:} the final approximation step --- since
$\Omega_\text{RF}\ll\Omega_\text{LO}$ is assumed,
$\alpha\kappa_\rho\Omega_\text{RF}\cos(\cdot)\ll1$, so
$e^{(\text{small})}\approx1+(\text{small})$. This reproduces Eq.~28
exactly, completing the derivation.

\subsection*{Eq.~(82) --- refined optimality condition}
\begin{equation}
\frac{\partial\kappa}{\partial\Omega_\text{LO}} = \frac{\partial(\bar P_0\kappa_\rho)}{\partial\Omega_\text{LO}} = \frac{\partial^2\bar P_0}{\partial\Omega_\text{LO}^2} = 0. \tag{82}
\end{equation}
\textbf{Significance:} a more refined optimality condition than
Eq.~79's, which only maximized $\kappa_\rho$ alone --- this one
maximizes the full intrinsic gain $\kappa=\alpha\bar P_0\kappa_\rho$
directly (i.e., accounts for how $\bar P_0$ itself changes with
$\Omega_\text{LO}$, not just $\kappa_\rho$).

\subsection*{Eq.~(83) --- refined optimal LO bias}
\begin{equation}
\Omega_\text{LO} = \sqrt{\bar A-1+\sqrt{\bar A^2-2\bar A+4}}\cdot\frac{\sqrt3}{3}\Gamma. \tag{83}
\end{equation}
\textbf{Significance:} solving Eq.~82 gives this $\bar A$-dependent
correction factor multiplying the simple Eq.~79 result. \textbf{Not yet
numerically checked in this project} against the simple
$\Omega_\text{LO,opt}=4.2294$~MHz from Eq.~79 --- worth computing if you
want to see whether it changes the match to the paper's stated 4.23~MHz
at all (likely a small correction, since $\bar A$ is typically not far
from a value that makes the correction factor close to 1, but this
should be verified numerically, not assumed).

\section{Appendix D: Fisher information / SNR derivation}

\subsection*{Eq.~(84) --- two-Lorentzian transmission model}
\begin{equation}
\mathcal T(\Delta_c;\Omega_\text{RF}) = 1-\mathcal C\sum_{i=\pm1}\frac{(\Gamma_\text{FWHM}/2)^2}{(\Delta_c-\frac{\Omega_\text{RF}}{2}i)^2+(\Gamma_\text{FWHM}/2)^2}, \qquad \mathcal C\in(0,1). \tag{84}
\end{equation}
\textbf{Terms:} $\mathcal C$ = a contrast parameter controlling how deep
the two Lorentzian dips are. $i=\pm1$ sums over the two AT-split peak
locations $\Delta_c=\pm\Omega_\text{RF}/2$.
\textbf{Significance:} \textbf{this is the exact model used for
Fig.~5} (``use QuTip toolkit to generate the EIT-AT spectrum as two
Lorentzian peaks''). \textbf{Important detail to verify if revisiting
Fig.~5}: this is written as a \emph{transmission-dip} model --- 1 minus
a sum of two Lorentzian \emph{peaks} in the dip/absorption sense --- not
a plain additive sum of two Lorentzian \emph{peaks} rising from a zero
baseline. If your own Fig.~5 implementation adds two Lorentzians
directly (peak-additive form) rather than subtracting them from a
constant background, double-check the two forms are actually
equivalent for your chosen normalization, since they are not
automatically the same thing outside of special cases.

\subsection*{Eq.~(85) --- Fisher information integral}
\begin{equation}
\mathcal I(\Omega_\text{RF}) = \frac{N_\text{ph}}{2\Gamma_\text{FWHM}}\int_{-\infty}^{+\infty}\frac{[\partial_{\mathcal R}\mathcal T(2\Delta_c/\Gamma_\text{FWHM};\mathcal R)]^2}{\mathcal T(2\Delta_c/\Gamma_\text{FWHM};\mathcal R)}\,d\left(\frac{2\Delta_c}{\Gamma_\text{FWHM}}\right). \tag{85}
\end{equation}
\textbf{Terms:} $N_\text{ph}$ = expected photon count within the
detection bandwidth (from a Poisson-distributed photon count model at
each detuning point).
\textbf{Significance:} the classical Fisher information for estimating
$\Omega_\text{RF}$ (via $\mathcal R$) from a full scan of the two-peak
transmission spectrum, treating each detuning-point measurement as
Poisson-distributed --- a rigorous information-theoretic starting point,
not a heuristic.

\subsection*{Eq.~(86) --- Fisher information bounds}
\begin{equation}
\mathcal I_0(\mathcal R) \le \mathcal I(\Omega_\text{RF}) \le \frac{1}{1-2\mathcal C}\mathcal I_0(\mathcal R). \tag{86}
\end{equation}
\textbf{Significance:} bounds the (generally intractable) exact
integral Eq.~85 between a computable closed form $\mathcal I_0(\mathcal
R)$ (Eq.~87) and a constant multiple of it --- lets the analysis
proceed without solving Eq.~85 exactly.

\subsection*{Eq.~(87) --- closed-form Fisher information lower bound}
\begin{equation}
\mathcal I_0(\mathcal R) = \frac{8N_\text{ph}\mathcal C^2}{\Gamma_\text{FWHM}^2}\frac{1}{1+\frac{1}{2\mathcal R^2}}. \tag{87}
\end{equation}
\textbf{Significance:} the explicit closed form of the lower bound
introduced in Eq.~86 --- already has exactly the same
$1/(1+\frac{1}{2\mathcal R^2})$ structure that will become $G(\mathcal
R)$ (Eq.~49/90).

\subsection*{Eq.~(88) --- exact Fisher information with contrast factor}
\begin{equation}
\mathcal I(\Omega_\text{RF}) = \frac{8N_\text{ph}\mathcal C^2}{\Gamma_\text{FWHM}^2}\frac{1}{1+\frac{1}{2\mathcal R^2}}\beta(\mathcal C), \qquad \beta(\mathcal C)\in\left[1,\tfrac{1}{1-2\mathcal C}\right]. \tag{88}
\end{equation}
\textbf{Terms:} $\beta(\mathcal C)$ = a contrast-only correction factor,
independent of $\mathcal R$.
\textbf{Significance:} the exact Fisher information, shown to differ
from the lower-bound Eq.~87 only by the $\mathcal R$-independent factor
$\beta(\mathcal C)$ --- for the paper's typical weak-probe simulations
($\mathcal C\lesssim0.2$), $\beta(\mathcal C)\le1.67$ ($\le2.22$~dB), a
small, absorbable correction.

\subsection*{Eq.~(89) --- Cram\'er--Rao bound}
\begin{equation}
\sigma_{\Omega_\text{RF}}^2 \ge \mathcal I^{-1}(\Omega_\text{RF}) = \frac{\Gamma_\text{FWHM}^2}{8N_\text{ph}\mathcal C^2}\left(1+\frac{1}{2\mathcal R^2}\right). \tag{89}
\end{equation}
\textbf{Significance:} the standard Cram\'er--Rao inequality applied to
Eq.~88 --- gives a lower bound on the variance of any unbiased
estimator of $\Omega_\text{RF}$ from the two-Lorentzian spectrum.

\subsection*{Eq.~(90) --- $G(\mathcal R)$, derived (not assumed)}
\begin{equation}
\frac{\Omega_\text{RF}^2}{\sigma_{\Omega_\text{RF}}^2} = \frac{\mathcal R^2}{1+\frac{1}{2\mathcal R^2}} \equiv G(\mathcal R). \tag{90}
\end{equation}
\textbf{Significance:} multiplying the classical link SNR term
$P_\text{Rx}|h|^2/\sigma_\text{Ry}^2$ by the RF-to-atomic conversion
gain $(\wp_\text{RF}/\hbar\Gamma_\text{FWHM})^2$ and this Fisher-derived
factor $G(\mathcal R)$ yields exactly Eq.~50 --- the whole point of
Appendix~D is to prove $G(\mathcal R)$ is a rigorous consequence of
information theory applied to AT-splitting estimation, not a
prescribed/arbitrary term.

\section{Full parameter table (Sec. V-A, as stated in the paper)}
\begin{center}
\begin{tabular}{@{}lll@{}}
\toprule
Symbol & Value & Meaning\\
\midrule
$L$ & 1~cm & vapor cell length\\
$N_0$ & $4.89\times10^{10}\,\text{cm}^{-3}=4.89\times10^{16}\,\text{m}^{-3}$ & total atomic density\\
States & $6S_{1/2}\to6P_{3/2}\to47D_{5/2}\to48P_{3/2}$ (Cs) & 4-level ladder\\
$\gamma_1$ & 0 & ground-state decay (none)\\
$\gamma_2/2\pi$ & 5.2~MHz & $|2\rangle$ decay rate\\
$\gamma_3/2\pi$ & 3.9~kHz & $|3\rangle=47D_{5/2}$ decay rate\\
$\gamma_4/2\pi$ & 1.7~kHz & $|4\rangle=48P_{3/2}$ decay rate\\
$\wp_{12}$ & $(2.5\,e a_0)^2$ & $|1\rangle$--$|2\rangle$ dipole moment (already squared)\\
$\wp_\text{RF}$ & $-1443.459\,ea_0$ & $|3\rangle$--$|4\rangle$ dipole moment\\
$f_\text{RF}$ & 6.9~GHz & RF carrier frequency\\
$\lambda_p$ & 852~nm & probe wavelength\\
$P_\text{in}$ & $20.7\,\mu$W & stated directly (inconsistent with Eq.~3 given $d$ below)\\
$d$ & 0.76~mm & probe 1/e$^2$ beam diameter\\
$\Omega_p/2\pi$ & 8~MHz & probe Rabi frequency\\
$\lambda_c$ & 510~nm & coupling wavelength\\
Coupling intensity & 17~mW & \\
$d_c$ & 1.95~mm & coupling 1/e$^2$ beam diameter\\
$\Omega_c/2\pi$ & 1~MHz & coupling Rabi frequency\\
$\epsilon_0$ & $8.854\times10^{-12}$~F/m & vacuum permittivity\\
$Z_0$ & 377~$\Omega$ & free-space impedance\\
$\Omega_\text{RF}/2\pi$ (LO-free) & 6~MHz & example RF field, LO-free figures\\
$\Omega_\text{RF}/2\pi$ (LO-dressed) & 20~kHz & example RF field, LO-dressed figures\\
$\Omega_\text{LO}/2\pi$ & 4.23~MHz & stated LO-dressed optimum\\
$B$ & 100~kHz & system bandwidth\\
$\Delta f$ & 15~kHz & LO/RF frequency difference\\
$G_\text{LNA}$ & 20~dB & LNA gain\\
$R_L$ & 50~$\Omega$ & load resistance\\
$D$ & 0.55~A/W & PD responsivity\\
$T$ & 290~K & temperature\\
$\eta_\text{eff}$ & 0.5 & PD quantum efficiency\\
$G_\text{Tx}$ & 2.15~dBi & Tx dipole antenna gain\\
$P_\text{Tx}$ & 30~dBm (1~W) & Tx power\\
$G_\text{Rx}$ (conventional only) & 2.15~dBi & Rx dipole antenna gain\\
\bottomrule
\end{tabular}
\end{center}

\section{Figure-by-figure equation chains}

\subsection*{Fig.~1, Fig.~2 --- schematic only}
No computation involved.

\subsection*{Fig.~3 --- $P_\text{out}(\Delta_c,\Omega_\text{RF})$, LO-free}
Eq.~6 ($\Omega=\Omega_\text{RF}$, sweep $\Delta_c,\Omega_\text{RF}$,
$\Delta_p=\Delta_\text{RF}=0$) $\to$ Eq.~7 $\to$ QuTiP steady state (NOT
Eq.~8) $\to$ Eq.~4 $\to$ Eq.~2. $P_\text{in}=20.7\,\mu$W used directly.
No Doppler averaging (Eq.~9--11 deliberately excluded).

\subsection*{Fig.~4 --- $P_\text{out}(\Delta_c,\mathcal R)$, LO-free}
Reuses Fig.~3's grid $\to$ Eq.~48 to relabel the axis as $\mathcal R$.
$\Gamma_\text{FWHM}$ choice unresolved analytically (see Eq.~77 caution
above); this project uses a measured value.

\subsection*{Fig.~5 --- $P_\text{out}(\Delta_c,d_\text{Tx-Rx})$ + SNR, LO-free}
Eq.~15 $\to$ Eq.~16/17 (needs $N_\text{atoms}$, not given) $\to$ Eq.~13
$\to$ Eq.~12 $\to$ Eq.~84 (two-Lorentzian, check dip-vs-peak form) $\to$
Eq.~50 for the SNR curve.

\subsection*{Fig.~6 --- LDR of LO-dressed receiver}
Eq.~6 ($\Omega=\Omega_\text{total}$) $\to$ QuTiP steady state $\to$
Eq.~2/4 $\to$ real $P_\text{out}(\Omega_\text{total})$. Overlay Eq.~78/79.
Shade LDR using Eq.~57 or Eq.~58 (BOTH need a chosen $\epsilon$).
Panel (b): $\Omega_\text{total}(t)$ phasor. Panel (c): interpolate
panel (a).

\subsection*{Fig.~7 --- received-power map, LO-dressed}
Same LDR machinery as Fig.~6, translated to $P_\text{Rx}$ via
Eq.~12/13, plotted against $\Omega_\text{LO}/2\pi$ (log axis, verified
directly from the real figure).

\subsection*{Fig.~8 --- SNR vs distance}
Eq.~14/15/16/17 $\to P_\text{Rx}$ $\to$ Eq.~50 / Eq.~52 (or 61) /
conventional SNR (using $\lambda_\text{RF}^2/4\pi$ as aperture).

\subsection*{Fig.~9 --- mutual information vs $P_\text{Rx}$}
SNR from the Fig.~8 chain $\to$ Eq.~59 (LO-free) or Eq.~60 (LO-dressed).

\subsection*{Fig.~10 --- SER vs $P_\text{Rx}$}
SNR from the Fig.~8 chain $\to$ standard $M$-QAM (conventional,
LO-dressed) or $M$-PAM (LO-free) SER formulas [59].

\section{Standing conventions already established in this project}
\begin{itemize}
\item $\Gamma_\text{FWHM}$ (LO-free, Fig.~4/5) = 1.2231~MHz, measured
from real QuTiP data, not an analytic candidate.
\item $\Gamma$ (LO-dressed, Fig.~6/7, Eq.~78) = 7.3255~MHz, analytic,
gives $\Omega_\text{LO,opt}=4.2294$~MHz (paper: 4.23~MHz).
\item $N_\text{atoms}$: never given anywhere in the paper or its three
cited references.
\item $\epsilon$ (Eq.~56--58): never given a numeric value. \textbf{Both}
Eq.~57 and Eq.~58 need it.
\item $P_\text{in}$/$d$ (20.7~$\mu$W / 0.76~mm) are mutually
inconsistent under Eq.~3 (would need 39.08~$\mu$W) --- a real,
version-confirmed paper-internal inconsistency, introduced in v3.
\item Real $N_0$ makes every QuTiP-computed curve sharper than the
paper's own plots (confirmed recurring across Fig.~3--6; a dedicated
audit ruled out Rydberg blockade, Eq.~2/4 form errors, and N$_0$/L
inconsistency as fixable explanations).
\end{itemize}

\end{document}
